\documentclass[11pt]{article}
\usepackage{amsmath,amssymb,amsthm}
\newtheorem{theorem}{Theorem}
\newtheorem{proposition}[theorem]{Proposition}
\title{A Three-Head Compact Reader for Five Independent Signs}
\author{}
\date{November 24, 2026}
\begin{document}
\maketitle

\paragraph{Status.}
This is a complete finite construction with an exact rational computer
audit of its 32 limiting sign words. It is not Lean-checked and is not
an asymptotic saving. Direct sums preserve the ratio $5/3$.

\section{The reader}
Let $u=(1,2,3,4,5)^{\mathsf T}$ and define
\[
 B_1=(2/3,8/3,-10/3,4/3,12)^{\mathsf T},\qquad
 B_2=(2/3,8/3,-1,85/6,29/3)^{\mathsf T},
\]
\[
 B_3=(-4/3,-16/3,13/3,-31/2,-65/3)^{\mathsf T}.
 \tag{1}
\]
Their sum is zero. Put
\[
 D_{ih}(e)=1-eu_i-e^4B_{ih},\qquad
 C^\pi_{ih}(e)=\frac1{D_{ih}(e)-e\pi(h)}, \tag{2}
\]
where $e>0$ and $\pi$ permutes $1,2,3$.

\begin{theorem}[Five independent strict signs]
There is $e_0>0$ such that for every $0<e<e_0$ and every
$y\in\{-1,1\}^5$, some permutation $\pi$ and some $z\in\mathbb R^3$
satisfy
\[
 y_i\sum_{h=1}^3 C^\pi_{ih}(e)z_h>0\qquad(1\leq i\leq5). \tag{3}
\]
Only three permutations are needed, one for each possible position
of token $2$.
\end{theorem}

\begin{proof}
Write $P_e(t)=\prod_{j=1}^3(t-ej)$ and
$w_j=1/\prod_{\ell\ne j}(j-\ell)$, so $w=(1/2,-1,1/2)$.
For $k=0,1,2$ define
\[
 a_h^{(k)}(e)=e^{-2-k}(1-e\pi(h))^k w_{\pi(h)}. \tag{4}
\]
For $\widehat C^\pi_{ih}=1/(1-eu_i-e\pi(h))$, partial fractions yield
\[
 \sum_h\widehat C^\pi_{ih}a_h^{(k)}
 =\frac{u_i^k}{P_e(1-eu_i)}. \tag{5}
\]
Indeed the proper rational function $(1-t)^k/(e^kP_e(t))$ has
residue $(1-ej)^kw_j/e^{2+k}$ at $t=ej$.
The exact difference
\[
 C^\pi_{ih}-\widehat C^\pi_{ih}
 =\frac{e^4B_{ih}}
 {(1-eu_i-e\pi(h)-e^4B_{ih})(1-eu_i-e\pi(h))} \tag{6}
\]
shows that $C^\pi(e)a^{(k)}(e)$ tends, for $k=0,1,2$, to
\[
 \mathbf1,\quad u,\quad u^{\circ2}+\sum_h w_{\pi(h)}B_h. \tag{7}
\]
If token $2$ occupies head $h$, the last vector is
$F'_h=u^{\circ2}-(3/2)B_h$. Its three possible values are
\[
 F'_1=(0,0,14,14,7),\quad
 F'_2=(0,0,21/2,-21/4,21/2),\quad
 F'_3=(3,12,5/2,157/4,115/2). \tag{8}
\]
The limiting planes therefore equal
$\operatorname{span}(\mathbf1,u,F_h)$, where
\[
 F_1=(0,0,-2,-2,-1),\quad
 F_2=(0,0,-2,1,-2),\quad
 F_3=(0,0,-2,1,2). \tag{9}
\]
In fact $F'_1=-7F_1$, $F'_2=-(21/4)F_2$, and
$F'_3=(37/4)F_3-6\mathbf1+9u$.

These three planes jointly meet every open orthant. Their exact
annihilator matrices are
\[
 G_1=\begin{pmatrix}1&-1&-1&1&0\\5&-6&-1&0&2\end{pmatrix},\quad
 G_2=\begin{pmatrix}5&-8&1&2&0\\2&-2&-1&0&1\end{pmatrix},
\]
\[
 G_3=\begin{pmatrix}5&-8&1&2&0\\4&-6&1&0&1\end{pmatrix}. \tag{10}
\]
Each has rank two and annihilates the corresponding three independent
columns. The script \texttt{audit\_three\_head\_compact\_reader.py}
checks all 32 words with exact rational arithmetic. For each word $y$
it finds a template and two columns of $G_h\operatorname{diag}(y)$
whose nonnegative linear combination cancels the sum of all its
columns. It thus constructs $v\in\ker G_h$ with $y_iv_i\geq1$,
solves for its rational coordinates in the three columns of the
plane, and verifies every strict sign. Template usage is 22, 6, and 4.

Use these fixed coordinates in the converging vectors (7), converting
between (8) and (9) by the displayed invertible affine changes.
Their limit has the required strict signs, hence so does the actual
vector for every sufficiently small positive $e$. Taking the minimum
of the 32 positive thresholds gives a common $e_0$.
\end{proof}

\section{Whole-cube admissibility}
\begin{proposition}[Independent columns]
For every $s>0$ and every $5\times s$ sign matrix, three original
admissible affine ratios give its strict signs on the one-hot row
by one-hot column grid. Their denominators have strictly negative
slopes and remain positive on the full $(5+s)$-dimensional cube.
\end{proposition}
\begin{proof}
Choose $e<e_0$ so small that every $p_{ih}=eu_i+e^4B_{ih}$ is positive
and $\sum_i p_{ih}+3se<1$ for every $h$. For column $j$, choose
$\pi_j,z_j$ by the theorem. On the whole cube set
\[
 R_h(x,v)=\frac{\sum_j z_{jh}v_j}
 {1-\sum_i p_{ih}x_i-e\sum_j\pi_j(h)v_j}. \tag{11}
\]
Every denominator is at least $1-\sum_i p_{ih}-3se>0$. On the
grid its summands are exactly those of (3).

To check the original atom conditions explicitly, any ratio
$(N_0+\sum_l n_lx_l)/(1-\sum_l a_lx_l)$ with $a_l>0$ and
$S=\sum_l a_l<1$ is an original atom. Choose $0<\alpha<1-S$,
$\rho_l=a_l/(1-\alpha)$, $\gamma=1-\sum_l\rho_l>0$,
$\delta=0$, $m_l=-n_l/a_l$, and
$\eta=N_0-\sum_l\rho_lm_l$. Substitution gives the stated numerator
and denominator exactly. Thus (11) changes no admissibility condition.
\end{proof}

\section{What this does not amplify}
For $H$ distinct tokens, use
$w_j=1/\prod_{\ell\ne j}(\lambda_j-\lambda_\ell)$ and a probe
perturbation $e^{2H-2}B_{ih}$. The same residue calculation gives
the ordinary limiting image
\[
 \operatorname{span}\{1,u,\ldots,u^{H-2},
          u^{H-1}+\sum_h w_{\pi(h)}B_h\}. \tag{12}
\]
For $H\geq2$, $\sum_jw_j=0$, so the union of these limiting planes
lies in a space of dimension at most $2H-1$. This rules out a
superconstant ratio for this limiting family itself. It does not
bound the finite-$e$ reader: query-dependent positive row rescalings
can reveal higher-order terms absent from an ordinary limit.
Nor does it bound different perturbation schemes.

An asymptotic improvement still requires a construction meeting all
orthants in dimension $\Omega(H\log H)$ with $O(H)$ heads and the
same whole-cube budget. The finite construction and the number of
permutation choices do not supply that construction.
\end{document}
